\section{The Constitutional Amendment Protocol}
\label{sec:amendments}

A constitution that cannot change is a prison; a constitution that
changes at the whim of a single entity is a dictatorship. The challenge
for Verifiable Institutions is to enable evolution without enabling
capture. We introduce the \emph{Constitutional Amendment Protocol},
which formally distinguishes between \emph{ordinary legislation}
(policy changes within the existing type structure) and
\emph{constituent acts} (structural changes to the type structure
itself). This distinction maps directly to the classical constitutional
law separation between \emph{pouvoir constitu\'{e}} (powers exercised
under the constitution) and \emph{pouvoir constituant} (the power to
rewrite the constitution)~\cite{sieyes1789}.

\subsection{Temporal Type Versioning}
\label{subsec:versioning}

The legal state of the system is not a static set of types but a
versioned chain of type definitions:
\[
  \mathcal{L} = \{\mathcal{L}_0, \mathcal{L}_1, \dots, \mathcal{L}_v\}
\]
where each $\mathcal{L}_v$ is a complete, self-consistent set of linear
resource type constraints, and transitions $\mathcal{L}_{v-1} \to
\mathcal{L}_v$ are governed by the amendment rules of
\S\ref{subsec:threshold}.

\begin{definition}[Versioned Validity]
\label{def:versioned-validity}
A verdict $V$ is \emph{constitutionally valid} only relative to a
specific legislative version $\mathcal{L}_v$. The validity predicate is:
\[
  \mathsf{Valid}_{\mathcal{L}_v}(V) \iff \mathsf{TypeCheck}(V,\; \Gamma_{\mathcal{L}_v}) = \top
\]
where $\Gamma_{\mathcal{L}_v}$ is the type context derived from
$\mathcal{L}_v$.
\end{definition}

To ensure backward compatibility, the Transparency Log $\Sigma$ stores
the hash $h_v = \mathsf{hash}(\mathcal{L}_v)$ at every block height
$b$ where a legislative transition occurs:
\[
  \Sigma \ni \langle b,\; h_v,\; t_v,\; \sigma_{\mathsf{amend}} \rangle
\]
An auditor $\mathcal{J}$ verifying a historical verdict $V_{b_0}$
retrieves $\mathcal{L}_v$ active at $b_0$ via the log and runs
$\mathsf{Valid}_{\mathcal{L}_v}(V_{b_0})$ using the type context of
\emph{that version}. This implements the constitutional principle of
\emph{nullum crimen sine lege}: decisions are judged by the law of
their time, not by retroactively applied future law. A court in 2028
auditing a 2024 decision does not apply 2028 type constraints.

\begin{remark}[Versioned Certificate Structure]
\label{rem:versioned-cert}
To make historical verification efficient without storing full type
contexts indefinitely, each compliance certificate embeds the
legislative version hash at issuance:
\[
  C_{v} = \langle \mathsf{hash}(\mathcal{L}_v),\;
                   \mathsf{hash}(E),\;
                   t_{\mathsf{issued}},\;
                   \sigma_{\mathcal{L}_v} \rangle
\]
An auditor holding $C_v$ and the log entry for $\mathcal{L}_v$ can
verify the certificate without re-executing the original binary,
requiring only the published legislative hash chain.
\end{remark}

\subsection{The Amendment Threshold: Checks on Regulatory Capture}
\label{subsec:threshold}

The principal risk in any amendment protocol is \emph{regulatory
capture}~\cite{stigler1971}: a compromised regulator colluding with the
operator to weaken the oversight guarantees encoded in $\mathcal{L}$.
If the Legislative Branch ($\mathcal{L}$) alone authorises amendments,
the system merely transfers absolutism from the operator to the
regulator---the same structural failure in a different location.

We prevent this by classifying amendments into two classes with
different authorisation requirements.

\begin{definition}[Ordinary Legislation vs.\ Constitutional Amendment]
\label{def:amendment-classes}
Let $\mathcal{L}^* \subset \mathcal{L}$ be the \emph{Stone Clauses}:
the subset of type definitions whose modification would alter the
capability partition $P_{\mathcal{L}} \cap P_{\mathcal{E}} \cap
P_{\mathcal{J}}$ established in Paper~5~\cite{btv-cc}.
\begin{itemize}
  \item $\Delta\mathcal{L}_{\mathsf{policy}}$: any change that modifies
  $\mathcal{L} \setminus \mathcal{L}^*$ (parameter updates, new
  evidence schema fields, threshold adjustments). Authorised by
  $\mathcal{L}$ alone: requires $\sigma_{\mathcal{L}}$.
  \item $\Delta\mathcal{L}^*$: any change that modifies a Stone
  Clause (the append-only property, the evidence binding invariant, the
  ZK-audit access). Requires \emph{Tripartite Ratification}
  (Definition~\ref{def:tripartite}).
\end{itemize}
\end{definition}

\begin{definition}[Tripartite Ratification]
\label{def:tripartite}
A constitutional amendment $\Delta\mathcal{L}^*$ is
\emph{constitutionally valid} if and only if it carries valid
signatures from representatives of all three independent powers:
\[
  \mathsf{Valid}(\Delta\mathcal{L}^*) \iff
  \sigma_{\mathcal{L}} \wedge \sigma_{\mathcal{J}} \wedge
  \sigma_{\mathcal{E}_{\mathsf{rep}}}
\]
where $\sigma_{\mathcal{J}}$ is the signature of an accredited Monitor
from the Judicial registry (Paper~5, \S5.2~\cite{btv-cc}), and
$\sigma_{\mathcal{E}_{\mathsf{rep}}}$ is the signature of an elected
operator representative from a published registry of certified
operators.
\end{definition}

\begin{theorem}[Amendment Soundness]
\label{thm:amendment-soundness}
Under Tripartite Ratification, no single principal can unilaterally
weaken the Constitutional Invariant $P_{\mathcal{L}} \cap
P_{\mathcal{E}} \cap P_{\mathcal{J}} = \emptyset$.
\end{theorem}

\begin{proof}
Suppose a regulator $\mathcal{L}$ is captured and wishes to amend
$\mathcal{L}^*$ to remove the append-only property (Type~2 Stone
Clause). This requires $\sigma_{\mathcal{J}}$. The Judicial branch
is structurally independent (Theorem~2, Paper~5~\cite{btv-cc}): its
signature key is not controlled by $\mathcal{L}$ or $\mathcal{E}$,
and its accreditation is managed by the public Monitor registry.
A captured regulator cannot forge $\sigma_{\mathcal{J}}$ without
compromising the ZK signing key of an accredited Monitor, which is
a separate cryptographic assumption. Symmetrically, an operator
cannot unilaterally weaken the evidence-binding constraint without
$\sigma_{\mathcal{L}} \wedge \sigma_{\mathcal{J}}$. Therefore, no
pairwise collusion of two principals suffices to amend
$\mathcal{L}^*$; all three must consent. \qed
\end{proof}

\begin{remark}[Relation to Supermajority Requirements]
Tripartite Ratification is the cryptographic analogue of the
supermajority and multi-chamber ratification requirements in
constitutional law (e.g., US Constitution Article~V: two-thirds of
Congress plus three-quarters of states; Brazilian CF/88 Art.~60:
two-thirds of both chambers in two sessions). The structural logic
is identical: important rules require broader consent than ordinary
legislation, making capture or hasty revision harder.
\end{remark}

\subsection{Sunset Clauses as Linear Resources}
\label{subsec:sunset}

Static laws accumulate technical debt. A compliance rule written for
2024 AI systems may be inadequate for 2030 systems without any actor
being at fault---society's understanding of algorithmic risk evolves.
Conventional legal systems rely on \emph{legislative inertia}: laws
remain in force until someone acts to repeal them. This creates a
default toward the status quo, which in practice means outdated
constraints persist long past their usefulness.

We invert this default by encoding legislative authority as a
\emph{linear resource with finite lifetime}.

\begin{definition}[Legislative Mandate Token]
\label{def:mandate}
A \emph{Legislative Mandate} is a linear resource:
\[
  M : \texttt{Mandate}[\mathcal{L}_v,\; t_{\mathsf{exp}}]
\]
where $t_{\mathsf{exp}}$ is the UTC timestamp of expiry,
determined at ratification time and recorded in $\Sigma$.
The \texttt{Mandate} token is \emph{linear}: it cannot be duplicated
or discarded, and it expires at $t_{\mathsf{exp}}$ regardless of
whether a renewal has been ratified.
\end{definition}

\begin{theorem}[Mandatory Renewal]
\label{thm:sunset}
In a Constitutional Code system with Mandate tokens, if no valid
$M'$ is ratified before $t_{\mathsf{exp}}$, the system enters
\emph{Constitutional Interregnum}: the type constructor
\texttt{Verdict::new} cannot be called and no new decisions
are produced until a renewed mandate is ratified.
\end{theorem}

\begin{proof}
The \texttt{Verdict::new} constructor has the signature:
\[
  \texttt{Verdict::new}(e: E,\; c: C,\; m: \&M) \to V
\]
where $m$ is a shared reference to a live mandate. At $t >
t_{\mathsf{exp}}$, the mandate $M$ is consumed by the runtime's
expiry check and cannot be borrowed. Any call to
\texttt{Verdict::new} yields a compile-time borrow error on $m$.
Since \texttt{DeliveryToken::seal} requires $V$ (Theorem~3,
Paper~5~\cite{btv-cc}), no decision can be delivered during
Interregnum. The system halts at the boundary between
$\mathcal{L}_v$ and $\mathcal{L}_{v+1}$ until Tripartite
Ratification produces $M'$. \qed
\end{proof}

\paragraph{The Inversion of Default.}
Conventional systems run under old rules until someone forces a
change---the inertia of the state. Constitutional Code runs under
no rules when the mandate expires, and restarts only when citizens
(via their three-power representatives) actively renew the
constitution. This is the inertia of the citizenry: the burden of
proof falls on those who want the system to operate, not on those
who want it reviewed. It encodes, in type theory, what Jeffersonians
argued in natural law: that each generation has the right to govern
itself under laws it has chosen~\cite{jefferson1789}.

\subsection{Composing the Amendment Protocol}
\label{subsec:composition}

The three mechanisms compose into a single coherent lifecycle:

\begin{enumerate}
  \item \textbf{Ordinary operation:} the system runs under
  $\mathcal{L}_v$ with mandate $M_v$ until $t < t_{\mathsf{exp}}$.
  \item \textbf{Policy update:} $\mathcal{L}$ proposes
  $\Delta\mathcal{L}_{\mathsf{policy}}$; signed by $\sigma_{\mathcal{L}}$;
  committed to $\Sigma$; new binary hash published; Monitors detect
  transition and verify backward compatibility of $C_v$ certificates.
  \item \textbf{Constitutional amendment:} any party proposes
  $\Delta\mathcal{L}^*$; requires $\sigma_{\mathcal{L}} \wedge
  \sigma_{\mathcal{J}} \wedge \sigma_{\mathcal{E}_{\mathsf{rep}}}$;
  committed to $\Sigma$ only after all three signatures are collected
  within a ratification window $[t_p, t_p + \tau]$.
  \item \textbf{Mandate renewal:} Tripartite process ratifies $M_{v+1}
  = \texttt{Mandate}[\mathcal{L}_{v+1}, t_{\mathsf{exp}}']$; system
  resumes from Interregnum.
  \item \textbf{Interregnum (failure to renew):} $t > t_{\mathsf{exp}}$
  without ratified $M'$; system halts; historical log $\Sigma$ remains
  publicly verifiable; past decisions remain auditable; new decisions
  are structurally impossible.
\end{enumerate}

Note that step~5 is not a bug but a constitutional feature: a system
that cannot be renewed is a system that should not run. The Interregnum
is the type-theoretic equivalent of an election: a structured pause
that forces deliberate recommitment before resuming power.
